\documentclass[11pt]{article}

\usepackage[margin=1in]{geometry}
\usepackage{amsmath,amssymb,mathtools}
\usepackage{booktabs,tabularx,array,longtable}
\usepackage{enumitem}
\usepackage{microtype}
\usepackage{graphicx}
\usepackage{pdflscape}
\usepackage[numbers,sort&compress]{natbib}
\usepackage[colorlinks=true,allcolors=blue]{hyperref}
\usepackage[nameinlink]{cleveref}
\usepackage{xcolor}
\usepackage{ragged2e}
\setlength{\emergencystretch}{3em}
\setlist[itemize]{leftmargin=1.5em}
\setlist[enumerate]{leftmargin=1.7em}

% Generated by scripts/build_dependency_analysis.py; do not edit.
\newcommand{\YWSImportClosure}{61}
\newcommand{\YWSDKImports}{18}
\newcommand{\YWSTauCetiImports}{18}
\newcommand{\DKUmbrellaImportClosure}{709}
\newcommand{\QuenchImportClosure}{80}
\newcommand{\YWSDKOverlap}{34}
\newcommand{\YWSQuenchOverlap}{0}


\newcommand{\Lean}{\textsc{Lean}}
\newcommand{\YWS}{Yu--Wang--Samworth}
\newcommand{\DK}{Davis--Kahan}
\newcommand{\TauCeti}{Tau Ceti}
\newcommand{\OurProject}{the present formalization project}

\title{Literature Review and Positioning Memo:\\
Formalizing the Davis--Kahan / Yu--Wang--Samworth\\
Perturbation Lineage}
\author{Internal working document for the formalization paper}
\date{Literature searched through 14 August 2026}

\begin{document}
\maketitle

\begin{abstract}
This document is deliberately broader than the related-work section intended for
\texttt{paper.tex}.  Its purpose is to record the current literature, identify
nearby claims that have already become crowded, and preserve candidate
positioning arguments before the paper is compressed for submission.  The main
conclusion is that the project should not be positioned merely as a formalization
of one convenient statistical form of the Davis--Kahan theorem, nor merely as a
large language model (LLM) autoformalization case study.  A more defensible
object is the \emph{Davis--Kahan / Yu--Wang--Samworth perturbation lineage}:
a deep formal audit of the 1970 Davis--Kahan source and its statistical
specialization, the foundational operator-theoretic infrastructure exposed by
that audit, source corrections including a finite-dimensional counterexample to
Davis--Kahan Proposition~4.4, and downstream reuse in contemporary statistical
formalization.  The human--LLM interaction and cost study is potentially novel
as an unusually instrumented record of such work, but should remain secondary
to the mathematics unless later analysis yields a particularly sharp empirical
result.
\end{abstract}

\tableofcontents
\clearpage

\section{Purpose, scope, and how to use this memo}
\label{sec:purpose}

This is an intentionally over-complete research memo.  It is not intended to be
copied wholesale into the eventual paper.  It serves four functions.

\begin{enumerate}
  \item It records the mathematical lineage around Davis--Kahan, Wedin, YWS,
        and later statistical perturbation theory, so that the paper's main
        contribution is located in its mathematical domain rather than only in
        the AI-formalization literature.
  \item It surveys the rapidly changing literature on large-scale and
        AI-assisted formalization through the current search date.  Several
        claims that would have sounded novel in early 2025---formalizing a whole
        paper, constructing missing library foundations, using blueprints,
        discovering source defects, or reporting token usage---now have direct
        precedents.
  \item It separates claims about \emph{formal proof capability},
        \emph{formalization workflow}, \emph{semantic/source fidelity}, and
        \emph{new mathematics}.  These are different contribution classes and
        should not be conflated.
  \item It records open comparison work that must be completed before a final
        novelty claim is made, especially a declaration-level comparison with
        the Davis--Kahan results recently added to the public Statistical
        Learning Theory in Lean repository.
\end{enumerate}

The review is current through 14 August 2026.  This date matters: large-scale
formalization is changing on a timescale of weeks, and some of the closest
papers discussed below appeared only days before this memo was written.

\subsection{Search procedure and confidence}
\label{sec:search-procedure}

This memo was assembled by a broad, non-systematic search followed by citation
snowballing.  Search surfaces included arXiv and publisher pages for primary
papers, public Lean repositories and project blueprints, and targeted web/code
searches.  Query families included the classical source names and theorem
terminology (Davis--Kahan, Yu--Wang--Samworth, Wedin, principal angles,
spectral subspaces, trace norm, unitarily invariant norms, Sylvester equations),
formalization terms (Lean, Coq, Isabelle, HOL, Mathlib, source fidelity,
autoformalization, blueprint, formal audit), and process terms (human--AI,
agents, prompts, tokens, cost, provenance, expert review).  Current 2026 papers
were checked against their primary arXiv records where possible rather than
relying on secondary summaries.

Two negative searches are worth recording but should not be treated as priority
proofs.  First, exact-string and theorem-number searches conducted on 14 August
2026 did not surface a discussion of the trace-norm failure of Davis--Kahan
Proposition~4.4 or an equivalent published counterexample.  Second, searches for
Yu--Wang--Samworth together with Lean/Coq/Isabelle/formalization terms did not
surface an obvious public full formalization of the paper.  Both conclusions are
fragile: theorem-number indexing is poor for older literature, code search does
not cover every repository or proof assistant, and an observation may appear
under different terminology.  The novelty claims remain open audit tasks in
\Cref{sec:open-search}.

\section{Executive positioning}
\label{sec:positioning}

\subsection{The safest high-level thesis}

A robust current framing is:

\begin{quote}
\emph{Source-faithful formalization of the Davis--Kahan perturbation lineage as
executable scholarship for mathematical statistics: verifying the cited
results, auditing the original sources, exposing and repairing incorrect
statements, and converting a large prerequisite burden into reusable formal
operator-theory and spectral-perturbation infrastructure.}
\end{quote}

This framing leaves open whether YWS or Davis--Kahan receives the most space in
the final paper.  It also makes the role of LLMs methodological rather than the
sole reason the mathematics is publishable.

\subsection{Why an exclusively YWS-centered story may be too narrow}

YWS is an excellent statistical endpoint: it is compact, widely useful, and its
population-gap formulation explains directly why statisticians often cite it
instead of the 1970 source \citep{yu2015davis,chen2021spectral}.  It is also a
clean object for a paper-keyed source census.  However, the repository history
and mathematical architecture indicate that the Davis--Kahan development is a
much larger object.  The hardest formal work includes principal-angle geometry,
unitarily invariant norms, direct rotations, spectral subspaces, Sylvester
machinery, compact/self-adjoint spectral infrastructure, and extensions across
finite- and infinite-dimensional settings.  The Proposition~4.4 counterexample
also belongs to the 1970 source, not to YWS.

The source-level dependency analysis in this repository supports this
qualification.  The YWS citation surface has a conservative local import
closure of \YWSImportClosure{} modules, including \YWSDKImports{} modules in the
DavisKahan tree and \YWSTauCetiImports{} modules in the \TauCeti{} staging tree.
By contrast, the Davis--Kahan umbrella closes over \DKUmbrellaImportClosure{}
local modules.  The two closures overlap in \YWSDKOverlap{} modules.
\Cref{tab:import-closures} gives the family-level counts.

% Generated by scripts/build_dependency_analysis.py; do not edit.
\begin{table}[tb]
\centering
\caption{Conservative source-level local import closures for three project surfaces. The closure includes the root module and is an upper bound on declaration-level proof dependence.}\label{tab:import-closures}
\small
\begin{tabular}{lrrrrrr}
\toprule
Surface & Total & YWS & DK & TauCeti & Quench & Other \\
\midrule
YWS citation surface & 61 & 25 & 18 & 18 & 0 & 0 \\
Davis--Kahan umbrella & 709 & 0 & 565 & 144 & 0 & 0 \\
Quench umbrella & 80 & 0 & 1 & 16 & 41 & 22 \\
\bottomrule
\end{tabular}
\end{table}


An import closure is an upper bound on declaration-level dependence: a module
may be imported without every declaration being used.  Nevertheless, the gross
asymmetry is informative.  YWS clearly reuses meaningful foundations developed
in the broader Davis--Kahan effort, but it does \emph{not} depend on the whole
Davis--Kahan development.  The eventual paper should therefore avoid the claim
that all of the deep foundational work was necessary merely to obtain YWS.  A
better statement is that the project formalized a broader perturbation-theory
lineage, within which YWS is a statistically important and source-auditable
endpoint.

\subsection{Claims that appear defensible, and claims that have become crowded}

The strongest prospective contribution bundle is the combination of:

\begin{itemize}
  \item a deep and source-faithful treatment of the classical Davis--Kahan
        development, together with a complete paper-level treatment of YWS;
  \item an explicit small finite-dimensional counterexample to the printed
        Davis--Kahan Proposition~4.4 and a repaired theorem surface;
  \item two documented printed defects in YWS, carefully distinguished from
        wholesale failure of its main theorem;
  \item substantial reusable foundations in spectral/operator theory and matrix
        analysis, with an explicit upstreaming path through \TauCeti{};
  \item verified downstream statistical reuse (the Quench theorem family);
  \item and, secondarily, an unusually detailed but explicitly incomplete record
        of model usage, resource cost, git provenance, and human intervention.
\end{itemize}

By contrast, current literature makes the following generic claims poor novelty
claims on their own: ``LLMs can write Lean,'' ``we formalized a whole paper or
textbook,'' ``the target needed missing Mathlib infrastructure,'' ``we used a
blueprint or dependency graph,'' ``formalization found source errors,'' ``we
preserved provenance,'' ``we report tokens/cost,'' or ``we performed a
source-faithful formalization.''  Each has substantial contemporary precedent;
several have multiple precedents in 2026 alone
\citep{milikic2026leanflow,nie2026formatheoria,zhang2026slt,wei2026asymptotic}.

\section{The mathematical perturbation lineage}
\label{sec:math-lineage}

\subsection{Classical operator and matrix perturbation theory}

The Davis--Kahan paper sits in a broader perturbation-theoretic tradition.
Kato's monograph supplies the classical operator-theoretic framework for
perturbations of spectra, spectral projections, and invariant subspaces
\citep{kato1976perturbation}.  Davis and Kahan's 1970 paper develops several
``sin $\Theta$,'' ``tan $\Theta$,'' and related rotation results for spectral
subspaces of self-adjoint operators \citep{davis1970rotation}.  Wedin's 1972
work gives the corresponding singular-subspace perturbation theory for the SVD
\citep{wedin1972perturbation}.  Stewart and Sun's later monograph organizes a
large body of matrix perturbation theory and became a standard reference for
numerical linear algebra \citep{stewart1990matrix}.

This lineage matters for positioning because ``the Davis--Kahan theorem'' is not
one theorem with one canonical statement.  There are multiple separation
conditions, norms, subspace representations, finite- versus infinite-dimensional
forms, and operator versus matrix surfaces.  Later work sharpened or reframed
parts of this theory.  Xie, for example, derived refinements of a Davis--Kahan
sin-$\Theta$ theorem using singular-value inequalities \citep{xie1997note};
Grubi\v{s}i\'c and Veseli\'c used weak Sylvester equations to derive spectral
subspace rotation bounds \citep{grubisic2005sylvester}; and Grubi\v{s}i\'c,
Kostrykin, Makarov, and Veseli\'c proved a tan-$2\Theta$ theorem for indefinite
quadratic forms \citep{grubisic2010tan}.  These are especially relevant to a
formal development whose dependency tree contains Sylvester-equation and
infinite-dimensional operator-theory branches.

The paper should therefore be precise about which Davis--Kahan statements it
formalizes, which are direct transcriptions of the 1970 source, which are
reusable modern interfaces, and which are strengthened or repaired results.  A
source census or theorem crosswalk is useful precisely because theorem names
such as ``Davis--Kahan'' are otherwise too coarse to support a fidelity claim.

\subsection{YWS as a statistical reformulation, not a replacement for the 1970 source}

Yu, Wang, and Samworth isolate a practical difficulty with direct statistical
use of Davis--Kahan: a common separation hypothesis is expressed using both
population and sample eigenvalues.  Their variant replaces this by a population
eigengap, gives explicit constants, and treats singular-vector subspaces for
rectangular matrices \citep{yu2015davis}.  In statistical applications this is
not a cosmetic restatement.  The population-only gap can be controlled at the
model level, whereas sample separation is itself random and may require a
second perturbation argument.

This is why YWS is a compelling paper-level formalization target.  It is small
enough that source fidelity can be stated concretely, but consequential enough
that subtle changes in hypotheses or constants matter.  It is also explicitly a
bridge from classical operator perturbation to modern statistics.  However,
that role should not erase the mathematically richer Davis--Kahan source from
which the present project grew.

\subsection{Modern statistical geometry around eigenspaces and singular subspaces}

Perturbation theory remains a central interface between random-matrix error
bounds and statistical conclusions.  The survey of Chen, Chi, Fan, and Ma
reviews spectral methods throughout high-dimensional statistics and data
science, with Davis--Kahan/Wedin-type bounds serving as standard tools for
turning matrix estimation error into eigenspace error \citep{chen2021spectral}.
Cape, Tang, and Priebe develop finer two-to-infinity norm geometry for singular
subspaces, motivated by high-dimensional statistical problems where Frobenius
or operator norm control can be too coarse \citep{cape2019two}.  Such work
illustrates why a formal library should expose the geometry of subspaces and
alignments as reusable mathematics, not hide it behind a single theorem
specialized to one norm.

For the final related-work section, one paragraph on this statistical lineage is
probably enough.  For the paper's mathematical introduction, however, it may be
worth showing explicitly how YWS sits between classical Davis--Kahan/Wedin and
later statistical spectral methods.

\section{Foundational formalization before the current LLM wave}
\label{sec:prellm}

Large formalization projects already established several practices that should
not be attributed to LLM agents.

\subsection{Scale and the cost of foundations}

Flyspeck's formal proof of the Kepler conjecture demonstrated that a major
mathematical result could require a large, carefully engineered formal corpus,
including auxiliary mathematics and substantial software infrastructure
\citep{hales2017kepler}.  The point relevant here is not simply scale: major
formalizations routinely uncover a long prerequisite chain that is invisible in
a concise informal proof.

The Liquid Tensor Experiment is even closer to the ``foundational debt'' story.
Scholze's challenge concerned a particular theorem in condensed mathematics
\citep{scholze2022liquid}, but the Lean project necessarily formalized a broad
supporting theory.  Its interactive blueprint linked mathematical exposition to
Lean declarations and exposed dependency graphs to track the project
\citep{liquidblueprint2023}.  The project also maintained a \texttt{for\_mathlib}
area for material intended to migrate upstream.  This is a direct historical
precedent for three aspects of the present project: source-linked tracking,
large prerequisite spillover, and separating project-local staging from reusable
library contribution.

The Lean Polynomial Freiman--Ruzsa (PFR) project likewise uses a blueprint and
dependency graph as a project-management and exposition layer
\citep{pfrblueprint2024}.  Thus a census, manifest, blueprint, or progress graph
is not intrinsically an AI contribution.  What may be empirically interesting
in our case is the interaction history by which tracking tools emerged and were
then adopted into human review---but the mechanism itself belongs to an
established formalization tradition.

\subsection{Library design is itself mathematical work}

Formalization papers frequently contribute interfaces and abstractions whose
value exceeds the original target.  Dedecker and Loreaux's formalization of the
continuous functional calculus emphasizes design decisions made for usability
and has been merged into Mathlib \citep{dedecker2025cfc}.  Marion's Ionescu--Tulcea
formalization similarly develops probability infrastructure that then supports
more general constructions \citep{marion2025ionescu}.  Recent matrix work
abstracts a common recursive schema across many matrix decompositions rather
than proving each decomposition in isolation \citep{ma2026matrixdecompositions}.

These examples strengthen, rather than weaken, the argument for reporting the
operator-theory foundations produced by the Davis--Kahan effort.  The right
claim is not that prerequisite spillover is unprecedented.  It is that the
formal audit identifies a concrete missing layer in spectral perturbation and
operator geometry, the project has already paid the construction cost, and the
resulting layer is being organized for reuse and upstream contribution.

\section{From theorem-proving benchmarks to LLM autoformalization}
\label{sec:llm-history}

\subsection{Benchmarks and local theorem proving}

Modern AI theorem proving developed initially around relatively local tasks.
MiniF2F provided a cross-system set of Olympiad-level formal statements
\citep{zheng2022minif2f}.  ProofNet coupled undergraduate-level natural-language
statements and proofs with Lean statements, explicitly targeting both
formalization and proving \citep{azerbayev2023proofnet}.  LeanDojo provided
programmatic access, extracted training data, a large theorem benchmark, and a
retrieval-augmented prover \citep{yang2023leandojo}.

Wu et al. showed that large language models could autoformalize a nontrivial
fraction of competition statements and use the resulting synthetic data to
improve theorem proving \citep{wu2022autoformalization}.  Draft, Sketch, and
Prove used informal proofs to guide formal proof search through generated
formal sketches \citep{jiang2023dsp}.  DeepSeek-Prover scaled synthetic Lean data
generation to millions of formal examples \citep{xin2024deepseek}.  A 2025
survey documents the rapidly expanding autoformalization landscape
\citep{weng2025survey}.

AlphaProof represents a different but equally important trajectory: formal
mathematics as a reinforcement-learning environment.  It trained on very large
numbers of autoformalized problems and reached medal-level combined performance
on the 2024 IMO with AlphaGeometry 2, demonstrating the power of kernel-checked
formal feedback at scale \citep{hubert2026alphaproof}.

These works make one point decisive for our paper: demonstrating that frontier
models can produce Lean code or close difficult proof obligations is no longer a
sufficient contribution.  The formalization itself, its fidelity, its
mathematical audit, its library contribution, or the empirical development
record has to carry the paper.

\subsection{Statement generation and proof generation are different problems}

A recurring lesson in the newer literature is that proving a fixed formal
statement and deciding what formal statement accurately represents a source are
distinct tasks.  A theorem prover can be perfectly sound about the wrong
translation.  This distinction becomes central for a project claiming to
formalize named results from Davis--Kahan and YWS rather than merely prove
related perturbation lemmas.

The eventual paper should exploit this distinction explicitly.  The source
census, theorem-number crosswalks, recorded corrections, and provenance review
address the translation boundary; the Lean kernel addresses proof validity
\emph{after} that boundary has been crossed.

\section{Document- and project-scale AI formalization in 2026}
\label{sec:scale2026}

This is the most rapidly changing related-work category and the strongest reason
not to make scale itself the headline claim.

\subsection{M2F, Automatic Textbook Formalization, and Atlas}

M2F presents an agentic two-stage framework for project-scale formalization,
first compiling a source into a coherent project with placeholder proofs and
then repairing the proofs under fixed signatures.  It reports 153,853 lines of
Lean generated from 479 pages of source material \citep{wang2026m2f}.  Automatic
Textbook Formalization reports a graduate textbook exceeding 500 pages,
approximately 130,000 lines, and 5,900 declarations produced by a very large
parallel population of agents \citep{gloeckle2026textbook}.  Formalizing
Mathematics at Scale generalizes this direction through AutoformBot and Atlas,
reporting 26 textbooks, more than 45,000 declarations, and about half a million
lines of Lean \citep{rammal2026scale}.

These projects establish that document-scale and textbook-scale formalization is
technically feasible with contemporary models.  Our paper should therefore not
compete on line count or number of agents.  A classical source audit in
operator/matrix perturbation theory is a different scientific object.

\subsection{LeanFlow: workflow ablations and token accounting}

LeanFlow is especially close to our prospective empirical process section.  It
formalizes two previously unformalized papers and experimentally varies model,
workflow, and tool access; importantly, it reports API calls and input/output
token usage \citep{milikic2026leanflow}.  Thus ``we counted tokens'' is already
an explicit research contribution in paper-to-Lean formalization.

Our dataset can still be complementary.  It spans a longer, organically evolving
project; includes multiple model families and interfaces; contains commit and
interaction traces; and has missingness mechanisms that themselves need to be
modeled.  But the comparison must be disciplined: token accounting conventions,
cache accounting, model pricing, and session structure differ.  A raw token
ratio between projects should not be presented as an efficiency ranking.

\subsection{LeanMarathon: long-horizon fidelity and an evolving blueprint}

LeanMarathon argues that long-horizon formalization fails through statement
drift, tangled dependencies, context decay, and nonlocal regressions, and uses
an evolving blueprint as both proof skeleton and shared system of record
\citep{zhang2026leanmarathon}.  This is a very close methodological precedent
for the project's emergent census/coordination tooling.  It also reinforces the
value of separating local theorem proving from long-horizon artifact
maintenance.

The paper can still report the census history if the interaction logs support an
interesting human--agent division of initiative.  It should not describe
``agents developed tracking tools'' as though project tracking itself were a
novel formalization method.

\subsection{Beyond the Library: constructing missing formal vocabulary}

Beyond the Library addresses research formalization in which the source uses
concepts absent from Mathlib.  Its agents dynamically extend the available
definitions and auxiliary lemmas before proving the target theorems, and the
system was applied to research papers from STOC \citep{moakhar2026beyondlibrary}.
This is the nearest generic precedent for the claim that research-level targets
force construction of missing formal foundations.

The distinction for our work should be mathematical and architectural: what
specific spectral/operator-theory substrate was missing, why it was required by
the Davis--Kahan source, how much is source-specific versus reusable, and what
portion is being upstreamed.  A dependency DAG and a carefully curated list of
foundational contributions would be more persuasive than a raw count of helper
lemmas.

\subsection{FormaTheoria: the closest current generic framing}

FormaTheoria, submitted 11 August 2026, is particularly important because its
problem statement overlaps several phrases we might otherwise have used as
headline novelty.  It describes large-scale Lean theory construction from
heterogeneous mathematical literature in terms of implicit dependency
discovery, correction of source defects, semantic fidelity, cross-source
misalignment, provenance preservation, structured review, and empirical process
analysis \citep{nie2026formatheoria}.

This is strong novelty pressure.  It means the final paper should not claim that
``formalization as executable scholarship'' is unique simply because the
workflow discovers dependencies and corrects sources.  Instead, the paper should
locate those general phenomena in a mathematically important statistical
perturbation lineage and emphasize the specific mathematical discoveries and
reusable theory produced.  If Proposition~4.4's trace-norm counterexample is new
to the mathematical literature, that fact is qualitatively different from the
generic observation that formalization can expose defects.

\section{Semantic fidelity, source audit, and counterexample-guided repair}
\label{sec:fidelity}

\subsection{Compilation is not semantic fidelity}

Beyond Compilation directly evaluates the gap between Lean compilation and
semantic faithfulness.  On its benchmark, the full tool-augmented agent achieves
much higher compile success than consensus semantic faithfulness, demonstrating
that a type-correct declaration can still omit hypotheses, alter domains, or
formalize a vacuous or otherwise incorrect claim \citep{zhang2026beyondcompilation}.
CriticLean similarly elevates semantic-fidelity criticism to a first-class part
of autoformalization \citep{peng2025criticlean}.  FormalScience studies semantic
drift in scientific formalization, including abstraction changes and notation
collapse \citep{meadows2026formalscience}.

These papers strongly support our insistence on source-level tracking.  A
paper-keyed census is useful not because it proves faithfulness by itself, but
because it creates reviewable obligations: theorem number, source statement,
formal declaration, status, correction/generalization flag, and provenance can
all be audited independently of compilation.

\subsection{Expert review after the sorries are gone}

Sorries Are Not the Hard Part reports an expert review of a semi-autonomous
formalization that compiled without placeholders but still had serious problems
with definitions, generality, file organization, and API design
\citep{ilin2026sorries}.  This is directly relevant to our paper's claims about
library quality and eventual \TauCeti{} contribution.  ``No sorry'' establishes
proof completion, not that the declarations form a good reusable mathematical
library.

A useful implication for our evaluation is that source fidelity and library
quality should be reported separately.  The YWS census can assess source
coverage.  A \TauCeti{} upstreaming/review process can assess API and library
quality.  Conflating them would overclaim both.

\subsection{Blueprint validation and synchronized source/formal metadata}

LeanArchitect generates and synchronizes blueprint information from annotations
in Lean declarations, infers dependency data, and exposes inconsistencies in
existing blueprints \citep{zhu2026leanarchitect}.  This provides another
precedent for compile-aware source/formal progress tooling.  It also suggests a
future improvement to our own census: over time, source identifiers and
provenance fields could migrate from paper-local scripts into declaration-level
metadata from which paper tables are generated automatically.

\subsection{Formal refutation and repaired statements}

MechGeo is a particularly relevant neighbor because it does not treat a failed
proof attempt as merely a search failure.  For historical geometry problems it
constructs Lean-checked counterexamples for statements that cannot be proved and
then proves expert-corrected versions \citep{shen2026mechgeo}.  This means the
generic pattern ``formalization exposed a false source statement and we repaired
it'' now has an explicit AI-formalization precedent.

The Davis--Kahan Proposition~4.4 result can nevertheless be a significant
mathematical contribution if the counterexample and repair are new to the
perturbation literature.  The paper should establish that novelty by searching
mathematical references and known errata, not by comparing only to formalization
papers.  It should also distinguish several levels of source problems:

\begin{itemize}
  \item a typographical/algebraic identity that is false as printed but locally
        repaired without changing the intended theorem;
  \item a boundary convention that makes a theorem specialization false or
        ill-posed but has a clear intended indexing repair;
  \item an invalid proof step with a correct theorem statement; and
  \item a genuinely false proposition requiring a counterexample and a changed
        theorem surface.
\end{itemize}

This taxonomy would prevent the YWS equation (4) typo, the YWS Theorem~3
boundary convention, the known Davis--Kahan Section~7 proof-line error, and
Proposition~4.4 from being rhetorically flattened into ``four false theorems.''

\section{Statistics and applied-science formalization}
\label{sec:stats}

This is the most important recent neighborhood for a math/statistics venue.

\subsection{Statistical Learning Theory in Lean 4}

Zhang, Lee, and Liu formalize empirical-process foundations for statistical
learning theory, including concentration and applications to regression.  They
explicitly report missing Mathlib infrastructure, human--AI collaboration, and
implicit assumptions exposed by formalization \citep{zhang2026slt}.  This is a
strong precedent for the claim that formalization of statistics forces
foundational mathematical work rather than merely encoding textbook statements.

The public repository has expanded beyond the February paper artifact.  Its July
2026 release advertises matrix spectral and perturbation theory including SVD,
Courant--Fischer, Eckart--Young--Mirsky, Weyl, and Davis--Kahan
\citep{zhang2026sltrepo}.  Therefore we should \textbf{not} claim ``the first
Lean formalization of Davis--Kahan'' without a much narrower qualifier.

A declaration-level comparison is still needed.  Questions to answer before
submission include:

\begin{enumerate}
  \item Which Davis--Kahan variant does the SLT repository state (eigenvector,
        projection, or general subspace; operator or Frobenius norm; which gap
        hypothesis)?
  \item Is its mathematical source the original 1970 paper, a textbook such as
        Vershynin, YWS, or another modern restatement?
  \item Does it cover the unitarily invariant norm/direct-rotation material,
        infinite-dimensional surfaces, or only the modern finite-dimensional
        statistical corollary?
  \item Does it make any source-fidelity claim about the 1970 paper or address
        Proposition~4.4?
\end{enumerate}

The likely outcome is that both projects contain valid Davis--Kahan results but
are formalizing different objects.  Our novelty language should be written only
after this comparison is complete.

\subsection{Source discipline in asymptotic statistics}

Wei et al. present a multi-agent formalization of asymptotic statistical theory
with an explicit ``hypothesis-disciplined'' audit: main-theorem assumptions must
be anchored to source prose, justified as encoding adapters, source-implied, or
rejected as unsupported strengthening \citep{wei2026asymptotic}.  This is very
close to the philosophy of a paper-keyed source census.  It further narrows the
novelty claim: source-faithful statistical formalization itself is now an active
subfield.

What remains distinctive is the particular mathematical lineage being audited,
the depth of the operator-theory prerequisite stack, and the specific source
corrections and downstream applications.

\subsection{Paper-organized applied mathematics libraries}

EconCSLib organizes Lean developments around research papers, retains
paper-facing statements, separates reusable shared infrastructure, and uses
post-formalization validation reports, dependency graphs, and human verification
of the translation boundary \citep{garg2026econcs}.  Again, paper-level
organization and dependency reports are no longer enough to distinguish a new
paper.

The Vlasov mean-field formalization is methodologically close in a different
way.  Miller describes the human as scoping definitions, steering decomposition,
and triaging library gaps while the AI executes the formal construction, and
explicitly distinguishes proving the statement as written from judging whether
that statement is the intended theorem \citep{miller2026vlasov}.  The work also
extracts a reusable optimal-transport layer from the application-specific
formalization.  This is a useful comparison for the relation between the
Davis--Kahan target and the reusable \TauCeti{} foundations.

Tensor-network formalization provides another close analogue from theoretical
physics: a multi-agent blueprint, periodic human review, extensive new
quantum/tensor infrastructure, and an explicit conclusion that preserving
mathematical intent is the main bottleneck \citep{lu2026tensor}.  Probability
formalization projects likewise emphasize the interface between source-facing
statements and Mathlib's more general abstractions
\citep{deng2026probability,marion2025ionescu}.

\subsection{Why a mathematical-statistics audience may still care}

The saturation of AI-formalization methodology does not imply saturation of
formalized statistics.  The relevant question for a mathematical-statistics
paper is whether a machine-checked development teaches us something concrete
about an important statistical tool.  A Davis--Kahan/YWS paper can answer yes if
it provides:

\begin{itemize}
  \item a trustworthy theorem map for variants statisticians actually cite;
  \item a correction to a classical source that affects a mathematically
        meaningful norm statement;
  \item reusable spectral-geometry interfaces for later statistical work;
  \item a demonstration that those interfaces shorten or clarify a downstream
        statistical formalization; and
  \item a quantitative account of what it cost to construct the missing formal
        layer, presented as supporting evidence rather than spectacle.
\end{itemize}

That is a domain contribution even if the underlying workflow uses techniques
that now have many AI-formalization precedents.

\section{Human--AI formalization process and resource studies}
\label{sec:process-literature}

\subsection{Controlled studies of formalization workflows}

Collins et al. study how people want to use AI for formalization and how they
actually use it in a controlled setting.  They report heterogeneous workflows
but a recurring desire to preserve high-level human control, and participants
frequently combine multiple AI tools \citep{collins2026humanai}.  This provides
an external framework for interpreting our prompt/interaction taxonomy: human
interventions are not merely failures of autonomy; they may be the mechanism by
which mathematical intent and project-level priorities remain under human
control.

\subsection{Case studies with public prompts and commits}

Ilin's Vlasov--Maxwell--Landau equilibrium formalization is one of the closest
comparators for transparent process reporting.  It reports ten days of work,
\$200 in cost, 229 human prompts, 213 git commits, zero human-written code, and
public process artifacts, together with failure modes including hypothesis creep
and definition alignment \citep{ilin2026vml}.  It therefore establishes a clear
precedent for reporting prompts, commits, cost, and human review in an
AI-assisted formalization.

Our process dataset differs in scale and messiness.  It spans multiple machines,
agents, models, model switches, chat surfaces, periods before instrumentation,
and measured sessions whose tokens cannot be defensibly assigned to individual
commits.  That messiness can itself be informative if modeled honestly, but it
means our results should not be compared to cleaner case studies as if every
``token'' or ``prompt'' were defined identically.

\subsection{Anthropic's Riemann-zeta case study}

Anthropic's August 2026 case study on Claude's mathematical capabilities is a
particularly salient comparison because it reports a long-running mathematical
agent workflow with released process information.  The report describes two
Claude Code sessions totaling 31 million output tokens, roughly 60 subagents,
and about 2,400 shell commands, with the human interaction characterized as
comparatively light \citep{anthropic2026zeta}.

Our likely contribution is not to beat those counts.  Rather, our records may
support a contrasting model of collaboration in which human interventions
include mathematical correction, theorem-scope decisions, source-fidelity
judgment, coordination, infrastructure constraints, and adoption/rejection of
agent-created management tools.  A useful empirical comparison would therefore
be about \emph{types of intervention and observability}, not only raw scale.

\subsection{Autonomy, interruption, and expertise in deployed coding agents}

Anthropic's large observational study of agent autonomy reports that experienced
Claude Code users approve actions more automatically while also interrupting
more often, shifting from action-by-action approval toward monitoring and
intervention \citep{mccain2026autonomy}.  Its later study of roughly 400,000
Claude Code sessions reports a division in which humans make more planning
choices while the agent makes more execution choices, and argues that domain
expertise continues to predict effective use \citep{anthropic2026expertise}.

These findings are unusually relevant to our interaction logs.  If our prompt
analysis finds that many human interventions are short corrections, priority
changes, missing-context injections, or coordination messages rather than
step-by-step Lean instructions, that pattern would fit a broader transition from
micro-management to supervisory steering.  The analysis should still be framed
as a case study, not as a causal result.

\subsection{Ephemeral interaction history as a provenance problem}

The present project's logs reveal a qualitative issue that is less prominent in
the current literature: the human participant does not reliably remember who
initiated every tool or project-management mechanism.  This is unsurprising in a
long-running agentic workflow.  Many instructions are operational, acted on
immediately, and never rehearsed as facts that the human expects to recall.

The source-census history gives a concrete example.  Retrospective memory
suggested that the agents had spontaneously created tracking machinery; the git
record shows that the human explicitly requested at least a manifest and
programmatic check, after which the agents elaborated the mechanism and the
human refined and incorporated it into review practice.  The contemporaneous
trace supports a more nuanced attribution than either memory alone or a simple
``human requested it'' label.

This motivates the phrase \emph{interaction archaeology}: logs can reconstruct
initiative, adoption, correction, and evolution that are otherwise ephemeral.
Whether this deserves space in the final paper depends on page budget and the
strength of the systematic evidence.  It is nevertheless worth preserving in
this memo because it affects how retrospective claims should be written:
contemporaneous traces should be preferred over unaided memory when available.

\subsection{Token accounting as an economic measurement problem}

Token usage is only meaningful economically when tied to the model and pricing
regime that consumed it.  LeanFlow already reports input/output tokens
\citep{milikic2026leanflow}; Anthropic's broader economic analysis also notes
that token use is a measure of compute whose interpretation can vary with task
and model \citep{anthropic2026expertise}.  Our ledger's per-model vectors allow
measured ordinary input, cache write, cache read, and output tokens to remain
separate.  That is preferable to a single project-wide token count or an average
cost per token.

The manuscript should report at least three layers separately:

\begin{enumerate}
  \item directly measured usage, by model and token class;
  \item measured but commit-unattributed session usage, which may be allocated
        only at coarser temporal/project resolution; and
  \item extrapolated usage for unmeasured work, with a missing-data model and
        uncertainty rather than a single reconstructed ``truth.''
\end{enumerate}

The limitation is scientifically important: work before telemetry and work in
chat interfaces may be absent not only from token counts but also from the
qualitative history of human intervention.  Missingness is therefore both a
resource-accounting problem and a provenance problem.

\section{Proof provenance, reuse, and upstreaming}
\label{sec:provenance}

\subsection{Formal proof reuse requires more than a bibliography}

The project records that proof ideas and implementation patterns were borrowed
from a former Spectra collaboration, external Lean libraries, and Mathlib pull
requests.  The eventual paper must cite these sources at the level needed to
make intellectual and code provenance clear.  A generic sentence saying that
``we used Mathlib'' is not enough when a particular proof architecture or block
of code was adapted from another formal development.

A useful audit taxonomy is:

\begin{enumerate}
  \item \textbf{close code adaptation}: code structure or proof terms closely
        follow an external formal source;
  \item \textbf{mathematical port}: the same formal mathematical strategy is
        ported but substantial API/type/universe rewriting makes the code new;
  \item \textbf{strategy donor}: an external proof suggested a route, but the
        final formal proof was independently re-derived in the local API;
  \item \textbf{ordinary library dependency}: the development invokes a public
        Mathlib theorem through its intended interface.
\end{enumerate}

The repository already has provenance blocks and external-source registries that
can seed this audit.  Before submission, the central theorems discussed in the
paper should each have a source/provenance check.  This is especially important
because LLM agents can reproduce patterns without the human remembering which
external code was inspected during the session.

\subsection{Upstreaming as evidence of reusable value}

A project-local proof of a famous theorem is useful, but a well-designed library
layer that survives upstream review has a different kind of value.  The Liquid
Tensor Experiment's \texttt{for\_mathlib} staging and recent formalization work
such as the continuous functional calculus provide precedents for treating
upstreamability as part of the engineering story
\citep{liquidblueprint2023,dedecker2025cfc}.

The current \TauCeti{} operator-theory roadmap can therefore be presented as an
ongoing reuse path rather than a novelty claim.  If major components are merged
before submission, the paper can report that as concrete evidence that the
foundational spillover is not merely auxiliary code needed to make one local
artifact compile.

\section{Closest-neighbor comparison}
\label{sec:neighbors}

\Cref{tab:neighbors} is intentionally qualitative.  Its purpose is to identify
which aspects of the project already have clear precedents and where the
combination may still be distinctive.  ``Source audit'' means an explicit
attempt to preserve/verify the claims of an informal source, not merely that a
paper motivated the formalization.  ``Process data'' means quantitative or
archived information about the construction process, not only a statement that
AI was used.

\begin{landscape}
\begin{table}[p]
\centering
\caption{Qualitative comparison with selected closest neighbors.  This is an
internal positioning table, not yet a submission-ready empirical claim.}
\label{tab:neighbors}
\small
\renewcommand{\arraystretch}{1.22}
\begin{tabularx}{\linewidth}{@{}p{3.0cm}p{6.2cm}p{6.2cm}X@{}}
\toprule
Work & What it establishes & Main overlap with the present project & Remaining distinction to investigate \\
\midrule
LeanFlow \citep{milikic2026leanflow}
 & Document-to-project formalization with source-fidelity concerns, workflow ablations, and API/token reporting.
 & Large source-grounded Lean development; missing project APIs; quantitative process/resource evidence.
 & Our strongest distinction must be mathematical audit and domain contribution, not token reporting or project scale. \\
FormaTheoria \citep{nie2026formatheoria}
 & Long-horizon theory construction with implicit-dependency discovery, semantic fidelity, provenance, source defects, and empirical process analysis.
 & Extremely close generic methodology and framing.
 & Perturbation-theory mathematics, the specific Davis--Kahan correction, complete YWS audit, and downstream statistical reuse must carry the distinction. \\
SLT in Lean \citep{zhang2026slt,zhang2026sltrepo}
 & Large statistical-learning formalization with foundational matrix/probability infrastructure; the public repository now includes a Davis--Kahan extension.
 & Closest statistics-domain neighbor and direct overlap in matrix perturbation.
 & Exact theorem-surface, source coverage, norm/generality, and provenance comparison is mandatory before any priority claim. \\
Asymptotic statistics \citep{wei2026asymptotic}
 & Source-grounded asymptotic statistics with explicit hypothesis discipline and multi-agent auditing.
 & Establishes source-faithful AI-assisted statistical formalization as an active category.
 & Our candidate distinction is the classical spectral-perturbation source audit, correction of a published statement, and operator-theory substrate. \\
Vlasov strategy game \citep{miller2026vlasov}
 & Long-horizon human-directed/AI-executed research formalization with reusable foundations and quantitative workflow observations.
 & Similar supervisory collaboration and foundational spillover.
 & Different mathematical domain (PDE/optimal transport); compare the types of human intervention rather than raw scale. \\
VML equilibrium \citep{ilin2026vml}
 & Transparent case study reporting prompts, commits, dollar cost, human review, and formalization failure modes.
 & Strong precedent for process/cost reporting tied to a mathematical artifact.
 & Our trace is larger and richer but substantially more incomplete; the value must come from observability and missing-data treatment, not cleaner accounting. \\
Tensor networks \citep{lu2026tensor}
 & Blueprint-driven, human-reviewed research formalization with substantial new library infrastructure and some new proof routes.
 & Close precedent for infrastructure growth and research-scale agent assistance.
 & Our source-audit/correction story and perturbation/statistics lineage remain the candidate mathematical differentiators. \\
Present project
 & Paper-keyed Davis--Kahan/YWS audit, extensive foundations, source repairs, downstream Quench reuse, and partial model-resolved process telemetry.
 & ---
 & Candidate novelty is the combination of a classical statistics/perturbation source audit, a potentially new Davis--Kahan correction, reusable operator-theory foundations, and verified downstream statistical use. \\
\bottomrule
\end{tabularx}
\end{table}
\end{landscape}

\section{Novelty pressure: claims to avoid, claims to investigate}
\label{sec:novelty-pressure}

\subsection{Claims we should not make without a narrow qualifier}

\begin{enumerate}
  \item \textbf{``The first Lean formalization of Davis--Kahan.''}
        The public SLT repository now advertises Davis--Kahan results
        \citep{zhang2026sltrepo}.  A narrower source/theorem/generality claim may
        be valid, but requires comparison.
  \item \textbf{``The first full-paper LLM formalization.''}
        M2F, LeanFlow, Beyond the Library, EconCSLib, and others already work at
        paper/project scale \citep{wang2026m2f,milikic2026leanflow,moakhar2026beyondlibrary,garg2026econcs}.
  \item \textbf{``The first project where formalization forced substantial new
        foundations.''}  This is a longstanding phenomenon and is explicit in
        recent AI-assisted work \citep{liquidblueprint2023,zhang2026slt,lu2026tensor}.
  \item \textbf{``The first use of a blueprint/census/dependency graph for
        formalization.''}  Liquid Tensor, PFR, LeanArchitect, LeanMarathon, and
        EconCSLib provide clear precedents
        \citep{liquidblueprint2023,pfrblueprint2024,zhu2026leanarchitect,zhang2026leanmarathon,garg2026econcs}.
  \item \textbf{``The first source-faithful AI formalization.''}
        LeanFlow, asymptotic statistics, FormaTheoria, and Beyond Compilation
        make fidelity an explicit object of study
        \citep{milikic2026leanflow,wei2026asymptotic,nie2026formatheoria,zhang2026beyondcompilation}.
  \item \textbf{``The first formalization to find a source error.''}
        Formalization has long exposed errors and missing assumptions; recent AI
        systems such as MechGeo and FormaTheoria explicitly center diagnosis and
        correction \citep{shen2026mechgeo,nie2026formatheoria}.
  \item \textbf{``The first AI formalization process study with token/cost
        reporting.''}  LeanFlow, VML, Anthropic's zeta case study, and broader
        agent-use studies are direct precedents
        \citep{milikic2026leanflow,ilin2026vml,anthropic2026zeta,mccain2026autonomy}.
\end{enumerate}

\subsection{Claims that may survive careful verification}

The following are plausible but should be treated as research questions until
checked.

\begin{enumerate}
  \item \textbf{A source-faithful formal treatment of the 1970 Davis--Kahan
        paper at unusual breadth.}  The exact scope should be described by a
        source census rather than a vague ``full paper'' label.
  \item \textbf{A new correction to Davis--Kahan Proposition~4.4.}  Search the
        perturbation literature, known errata, citation graph, and later papers
        that invoke the proposition.  Establish whether the trace-norm
        counterexample or equivalent observation has appeared previously.
  \item \textbf{A source-faithful complete YWS formalization including its
        sharpness examples and source repairs.}  We have not yet found another
        public full YWS formalization, but absence from search results is not a
        proof of priority.
  \item \textbf{An unusually broad reusable formal spectral-perturbation layer
        motivated by one classical source.}  Quantify breadth through a theorem
        crosswalk and dependency architecture, not raw line count.
  \item \textbf{An integrated audit-to-reuse story:} classical source
        $\rightarrow$ correction $\rightarrow$ modern statistical variant
        $\rightarrow$ contemporary downstream statistical theorem.  The
        combination may be more distinctive than any one step.
  \item \textbf{A uniquely granular naturalistic interaction dataset for
        research-mathematics formalization.}  Compare the actual fields and
        missingness mechanisms to published prompt/agent studies before making a
        priority claim.
\end{enumerate}

\section{Recommended eventual related-work cut}
\label{sec:recommended-cut}

If the final venue affords only roughly one page of related work, the material
above can be compressed into five paragraphs.

\paragraph{1. Perturbation theory and statistics.}
Davis--Kahan, Wedin, YWS, and one modern statistical spectral reference.  This
establishes the mathematical lineage and explains why YWS is statistically
important without implying that it subsumes the full 1970 source
\citep{davis1970rotation,wedin1972perturbation,yu2015davis,chen2021spectral}.

\paragraph{2. Large formalization and foundational spillover.}
One pre-LLM precedent (Liquid Tensor or Flyspeck) plus one operator/library paper
such as the continuous functional calculus.  The purpose is to acknowledge that
large prerequisite layers, dependency tracking, and upstreaming are established
formalization practice \citep{hales2017kepler,liquidblueprint2023,dedecker2025cfc}.

\paragraph{3. Project-scale AI formalization.}
Cite a small representative set: M2F/Atlas for scale, LeanFlow for workflow and
tokens, and FormaTheoria for the closest source-audit/provenance framing
\citep{wang2026m2f,rammal2026scale,milikic2026leanflow,nie2026formatheoria}.
Do not enumerate every 2026 system in the main paper.

\paragraph{4. Semantic fidelity and statistics.}
Use Beyond Compilation plus the two statistics neighbors: SLT in Lean and
hypothesis-disciplined asymptotic statistics
\citep{zhang2026beyondcompilation,zhang2026slt,wei2026asymptotic}.  This both
motivates our source census and acknowledges that source fidelity in statistics
is no longer unique.

\paragraph{5. Human--AI process.}
If the empirical section survives: LeanFlow, VML, Collins et al., and Anthropic's
zeta/autonomy work are enough to locate token accounting and interaction analysis
\citep{milikic2026leanflow,ilin2026vml,collins2026humanai,anthropic2026zeta,mccain2026autonomy}.
The appendix can hold the broader process literature.

The main paper should spend the citations saved here on the specific mathematical
literature needed to establish the novelty of the Proposition~4.4 correction and
the exact Davis--Kahan theorem surfaces.

\section{Open searches and comparison work before submission}
\label{sec:open-search}

This review is deep enough for positioning but not yet the final systematic
literature audit.  The following tasks remain.

\begin{enumerate}
  \item \textbf{Davis--Kahan Proposition~4.4 citation audit.}  Search by
        proposition number, distinctive theorem language, direct-rotation
        terminology, Fan--Hoffman/trace-norm formulations, and later errata.
        Trace citations backward and forward from papers that use Section~4 of
        Davis--Kahan.  This is the highest-priority mathematical novelty check.
  \item \textbf{Exact comparison to the SLT repository.}  Record theorem
        signatures, source references, norm/gap assumptions, and dependencies.
        A public repository claiming Davis--Kahan makes a vague priority claim
        untenable, but it may be mathematically far narrower than our source
        development.
  \item \textbf{Search other proof assistants.}  Current web searches did not
        reveal an obvious Coq or Isabelle formalization of Davis--Kahan/YWS, but
        negative keyword search is weak evidence.  Search AFP, Coq packages,
        HOL Light corpora, and theorem-name variants directly.
  \item \textbf{YWS priority search.}  Search GitHub/code indexes, Mathlib PRs,
        Lean projects, and supplementary artifacts for Yu--Wang--Samworth by
        theorem content as well as title/authors.
  \item \textbf{Perturbation-theory formalization beyond Davis--Kahan.}
        Search for formal Weyl, Courant--Fischer, principal angles, Wedin,
        Stewart--Sun, Kato spectral projections, Riccati/subspace rotation, and
        unitarily invariant norm results.  These determine which foundational
        contributions are genuinely new to formal libraries.
  \item \textbf{Proof-provenance citations.}  Complete the declaration-level
        audit of Spectra-derived proofs, Mathlib PR strategies, and external Lean
        code before writing the implementation section.
  \item \textbf{Bibliometrics.}  If the introduction says Davis--Kahan or YWS
        is ``highly cited,'' use a reproducible bibliometric snapshot and date.
        Citation counts are unstable and should not be copied from a search-engine
        sidebar without provenance.
  \item \textbf{Venue-specific review.}  Before submission, repeat the search
        for papers posted after 14 August 2026 and for accepted MathStat/AI-for-Math
        papers not yet indexed by arXiv search.
\end{enumerate}

\section{Working conclusion}
\label{sec:conclusion}

The literature does not suggest that this project has become unpublishable
because AI formalization is exploding.  It suggests that the paper has to be
more mathematically specific than it would have needed to be a year earlier.
Scale, agentic workflow, source fidelity, gap discovery, foundational spillover,
and process telemetry are now all active research themes.  None should be sold
as the sole novelty.

The project is strongest when treated as a formal mathematical investigation of
a classical perturbation-theory lineage with direct statistical importance.  In
that framing, the hard Davis--Kahan work matters on its own; YWS supplies a
source-faithful statistical endpoint rather than the only reason the
foundations exist; the Proposition~4.4 counterexample can stand as a potentially
new mathematical correction; the \TauCeti{} work demonstrates reusable
spillover; and Quench demonstrates downstream statistical use.  The interaction
and accounting study then answers a secondary but genuinely interesting
question: what did it take, in practice, for a human and contemporary LLM agents
to construct and audit this much formal mathematics?

That combination remains meaningfully different from benchmark theorem proving,
textbook-scale translation, and generic agentic formalization frameworks.  The
next novelty-critical step is not more AI-literature collection; it is a careful
mathematical citation audit of Proposition~4.4 and a theorem-by-theorem
comparison with the newly public Davis--Kahan material in adjacent Lean
projects.

\clearpage
\bibliographystyle{plainnat}
\bibliography{references}

\end{document}
